%% kapa: exam=1 round=37 year=2026 period=1 form=A subject=civil_law subject_ko=민법 seq=3 number=3
%% source: Q-Net 감정평가사 기출문제 (cst00309) artlSeq=5258924 file=raw/1008/5258924/2026년 제37회 감정평가사 1차 1교시.pdf
%% conversion: kordoc 4.24.51 -> tools/gen1.py (auto); review=proofread-sonnet; flags=none
\begin{question}{3}
성년후견에 관한 설명으로 옳지 않은 것은?
\begin{choices}
  \item 가정법원은 취소할 수 없는 피성년후견인의 법률행위의 범위를 정할 수 있다.
  \item 검사 또는 지방자치단체의 장도 특정후견심판의 청구권자에 포함된다.
  \item 가정법원은 특정후견의 심판을 하는 경우, 특정후견의 기간 또는 사무의 범위를 정하여야 한다.
  \item 가정법원은 본인의 복리상 필요한 경우에는 본인의 의사에 반하여 특정후견개시의 심판을 할 수 있다.
  \item 가정법원이 피한정후견인에 대하여 성년후견개시의 심판을 할 때에는 종전의 한정후견의 종료 심판을 한다.
\end{choices}
\end{question}
\answer{4}
